% How Small Can You Go? -- first draft.
% Build: cd paper && tectonic main.tex   (or latexmk -pdf main.tex)
% Numbers in generated/ come from scripts/make_paper_tables.py; do not edit them by hand.
%
% Class: IEEEtran journal mode (arXiv version; ports to the IEEE Access template, which
% is IEEEtran-based). For MDPI Sensors, move the text into the MDPI template; section
% order already follows Introduction / Related Work / Materials and Methods / Results /
% Discussion / Conclusions.
\documentclass[journal]{IEEEtran}

\usepackage[T1]{fontenc}
\usepackage{amsmath,amssymb}
\usepackage{graphicx}
\usepackage{booktabs}
\usepackage{multirow}
\usepackage{xcolor}
\usepackage{url}
\usepackage[hidelinks]{hyperref}
\usepackage{cite}

% ---- Draft markers ------------------------------------------------------------
% \pending{...}: result or detail that depends on experiments not yet run.
% \authornote{...}: something only the author can supply or decide.
% Set \draftfalse before submission; then any remaining marker breaks the build.
\newif\ifdraft \drafttrue
\ifdraft
  \newcommand{\pending}[1]{{\color{red!75!black}[\textbf{pending:} #1]}}
  \newcommand{\authornote}[1]{{\color{blue!70!black}[\textbf{author:} #1]}}
\else
  \newcommand{\pending}[1]{\errmessage{Unresolved pending marker}}
  \newcommand{\authornote}[1]{\errmessage{Unresolved author note}}
\fi

% Generated by scripts/make_paper_tables.py -- do not edit.
\newcommand{\AblLargeCpuDefaultFpGiB}{3.72}
\newcommand{\AblLargeCpuDefaultLoadS}{2.7}
\newcommand{\AblLargeCpuDefaultMedianMs}{5071}
\newcommand{\AblLargeCpuDefaultRssGiB}{6.42}
\newcommand{\AblLargeCpuDefaultSysGiB}{4.71}
\newcommand{\AblLargeCpuNoBothFpGiB}{3.86}
\newcommand{\AblLargeCpuNoBothLoadS}{1.0}
\newcommand{\AblLargeCpuNoBothMedianMs}{6161}
\newcommand{\AblLargeCpuNoBothRssGiB}{3.92}
\newcommand{\AblLargeCpuNoBothSysGiB}{5.28}
\newcommand{\AblLargeCpuNoMmapFpGiB}{4.24}
\newcommand{\AblLargeCpuNoMmapLoadS}{1.7}
\newcommand{\AblLargeCpuNoMmapMedianMs}{5471}
\newcommand{\AblLargeCpuNoMmapRssGiB}{4.29}
\newcommand{\AblLargeCpuNoMmapSysGiB}{2.05}
\newcommand{\AblLargeCpuNoRepackFpGiB}{1.20}
\newcommand{\AblLargeCpuNoRepackLoadS}{0.8}
\newcommand{\AblLargeCpuNoRepackMedianMs}{6270}
\newcommand{\AblLargeCpuNoRepackRssGiB}{3.92}
\newcommand{\AblLargeCpuNoRepackSysGiB}{1.36}
\newcommand{\AblLargeGpuDefaultFpGiB}{1.09}
\newcommand{\AblLargeGpuDefaultLoadS}{0.8}
\newcommand{\AblLargeGpuDefaultMedianMs}{576}
\newcommand{\AblLargeGpuDefaultRssGiB}{3.72}
\newcommand{\AblLargeGpuDefaultSysGiB}{4.54}
\newcommand{\AblLargeGpuNoMmapFpGiB}{4.27}
\newcommand{\AblLargeGpuNoMmapLoadS}{1.2}
\newcommand{\AblLargeGpuNoMmapMedianMs}{560}
\newcommand{\AblLargeGpuNoMmapRssGiB}{4.25}
\newcommand{\AblLargeGpuNoMmapSysGiB}{2.09}
\newcommand{\AblSmallCpuDefaultFpGiB}{0.71}
\newcommand{\AblSmallCpuDefaultLoadS}{0.4}
\newcommand{\AblSmallCpuDefaultMedianMs}{1020}
\newcommand{\AblSmallCpuDefaultRssGiB}{1.09}
\newcommand{\AblSmallCpuDefaultSysGiB}{0.58}
\newcommand{\AblSmallCpuNoBothFpGiB}{0.73}
\newcommand{\AblSmallCpuNoBothLoadS}{0.4}
\newcommand{\AblSmallCpuNoBothMedianMs}{1365}
\newcommand{\AblSmallCpuNoBothRssGiB}{0.76}
\newcommand{\AblSmallCpuNoBothSysGiB}{0.62}
\newcommand{\AblSmallCpuNoMmapFpGiB}{0.84}
\newcommand{\AblSmallCpuNoMmapLoadS}{0.4}
\newcommand{\AblSmallCpuNoMmapMedianMs}{1125}
\newcommand{\AblSmallCpuNoMmapRssGiB}{0.87}
\newcommand{\AblSmallCpuNoMmapSysGiB}{0.79}
\newcommand{\AblSmallCpuNoRepackFpGiB}{0.38}
\newcommand{\AblSmallCpuNoRepackLoadS}{0.4}
\newcommand{\AblSmallCpuNoRepackMedianMs}{1295}
\newcommand{\AblSmallCpuNoRepackRssGiB}{0.77}
\newcommand{\AblSmallCpuNoRepackSysGiB}{0.83}
\newcommand{\AblSmallGpuDefaultFpGiB}{0.40}
\newcommand{\AblSmallGpuDefaultLoadS}{0.4}
\newcommand{\AblSmallGpuDefaultMedianMs}{119}
\newcommand{\AblSmallGpuDefaultRssGiB}{0.70}
\newcommand{\AblSmallGpuDefaultSysGiB}{0.85}
\newcommand{\AblSmallGpuNoMmapFpGiB}{0.89}
\newcommand{\AblSmallGpuNoMmapLoadS}{0.4}
\newcommand{\AblSmallGpuNoMmapMedianMs}{120}
\newcommand{\AblSmallGpuNoMmapRssGiB}{0.84}
\newcommand{\AblSmallGpuNoMmapSysGiB}{0.66}
\newcommand{\BenchDevDown}{63}
\newcommand{\BenchDevHard}{92}
\newcommand{\BenchDevNormal}{62}
\newcommand{\BenchNCand}{777}
\newcommand{\BenchNDev}{217}
\newcommand{\BenchNKept}{711}
\newcommand{\BenchNMcfScen}{24}
\newcommand{\BenchNRejected}{66}
\newcommand{\BenchNRelabelled}{40}
\newcommand{\BenchNSeqDev}{24}
\newcommand{\BenchNSeqTest}{69}
\newcommand{\BenchNTest}{494}
\newcommand{\BenchNUrSeq}{69}
\newcommand{\BenchTestDown}{170}
\newcommand{\BenchTestHard}{171}
\newcommand{\BenchTestNormal}{153}
\newcommand{\PilotCpu}{Apple M4 Pro}
\newcommand{\PilotCpuGpuAgree}{210/210}
\newcommand{\PilotFloorHits}{81/81}
\newcommand{\PilotLargeAAuroc}{1.00}
\newcommand{\PilotLargeAFaHard}{2/10}
\newcommand{\PilotLargeAFaHardCI}{0--50\%}
\newcommand{\PilotLargeAFaNormal}{0/10}
\newcommand{\PilotLargeAFone}{0.91}
\newcommand{\PilotLargeBAuroc}{1.00}
\newcommand{\PilotLargeBFaHard}{6/10}
\newcommand{\PilotLargeBFaNormal}{0/10}
\newcommand{\PilotLargeBFone}{0.77}
\newcommand{\PilotLargeBtwoAuroc}{1.00}
\newcommand{\PilotLargeBtwoFaHard}{3/10}
\newcommand{\PilotLargeBtwoFaNormal}{0/10}
\newcommand{\PilotLargeBtwoFone}{0.87}
\newcommand{\PilotLargeCpuDecodeMs}{33}
\newcommand{\PilotLargeCpuFootprintGiB}{3.7}
\newcommand{\PilotLargeCpuMedianMs}{4270}
\newcommand{\PilotLargeCpuPrefillPct}{99}
\newcommand{\PilotLargeCpuRssGiB}{6.4}
\newcommand{\PilotLargeCpuStreams}{0}
\newcommand{\PilotLargeGpuDecodeMs}{28}
\newcommand{\PilotLargeGpuFootprintGiB}{1.1}
\newcommand{\PilotLargeGpuMedianMs}{530}
\newcommand{\PilotLargeGpuPrefillPct}{92}
\newcommand{\PilotLargeGpuRssGiB}{3.7}
\newcommand{\PilotLargeGpuStreams}{1}
\newcommand{\PilotLlamaBuild}{10090}
\newcommand{\PilotNAnswers}{621}
\newcommand{\PilotNLenient}{0}
\newcommand{\PilotNUnparsed}{0}
\newcommand{\PilotOs}{26.5.2}
\newcommand{\PilotQeightBfAgree}{59/60}
\newcommand{\PilotQwenAAuroc}{1.00}
\newcommand{\PilotQwenAFaHard}{2/10}
\newcommand{\PilotQwenAFaNormal}{0/10}
\newcommand{\PilotQwenAFone}{0.91}
\newcommand{\PilotQwenBAuroc}{1.00}
\newcommand{\PilotQwenBFaHard}{4/10}
\newcommand{\PilotQwenBFaNormal}{0/10}
\newcommand{\PilotQwenBFone}{0.83}
\newcommand{\PilotQwenBtwoAuroc}{1.00}
\newcommand{\PilotQwenBtwoFaHard}{0/10}
\newcommand{\PilotQwenBtwoFaNormal}{0/10}
\newcommand{\PilotQwenBtwoFone}{1.00}
\newcommand{\PilotQwenCpuDecodeMs}{18}
\newcommand{\PilotQwenCpuFootprintGiB}{2.8}
\newcommand{\PilotQwenCpuMedianMs}{4967}
\newcommand{\PilotQwenCpuPrefillPct}{99}
\newcommand{\PilotQwenCpuRssGiB}{4.8}
\newcommand{\PilotQwenCpuStreams}{0}
\newcommand{\PilotQwenGpuDecodeMs}{15}
\newcommand{\PilotQwenGpuFootprintGiB}{1.0}
\newcommand{\PilotQwenGpuMedianMs}{523}
\newcommand{\PilotQwenGpuPrefillPct}{92}
\newcommand{\PilotQwenGpuRssGiB}{2.8}
\newcommand{\PilotQwenGpuStreams}{1}
\newcommand{\PilotRamGB}{48}
\newcommand{\PilotRecallWilsonCI}{72--100\%}
\newcommand{\PilotSameRoomNormalRejectA}{9/9}
\newcommand{\PilotSmallAAuroc}{1.00}
\newcommand{\PilotSmallAFaHard}{5/10}
\newcommand{\PilotSmallAFaNormal}{0/10}
\newcommand{\PilotSmallAFone}{0.80}
\newcommand{\PilotSmallBAuroc}{0.89}
\newcommand{\PilotSmallBFaHard}{9/10}
\newcommand{\PilotSmallBFaNormal}{10/10}
\newcommand{\PilotSmallBFone}{0.51}
\newcommand{\PilotSmallBtwoAuroc}{0.94}
\newcommand{\PilotSmallBtwoFaHard}{9/10}
\newcommand{\PilotSmallBtwoFaNormal}{9/10}
\newcommand{\PilotSmallBtwoFone}{0.53}
\newcommand{\PilotSmallCpuDecodeMs}{5}
\newcommand{\PilotSmallCpuFootprintGiB}{0.7}
\newcommand{\PilotSmallCpuMedianMs}{916}
\newcommand{\PilotSmallCpuPrefillPct}{97}
\newcommand{\PilotSmallCpuRssGiB}{1.1}
\newcommand{\PilotSmallCpuStreams}{1}
\newcommand{\PilotSmallGpuDecodeMs}{6}
\newcommand{\PilotSmallGpuFootprintGiB}{0.4}
\newcommand{\PilotSmallGpuMedianMs}{112}
\newcommand{\PilotSmallGpuPrefillPct}{76}
\newcommand{\PilotSmallGpuRssGiB}{0.7}
\newcommand{\PilotSmallGpuStreams}{8}


\newcommand{\pd}{\textit{person-down}}
\newcommand{\promptA}{\texttt{A}}
\newcommand{\promptB}{\texttt{B}}
\newcommand{\promptBtwo}{\texttt{B2}}

\begin{document}

\title{How Small Can You Go? Accuracy--Efficiency Trade-offs of Edge Vision--Language Models for Person-Down Detection in Industrial Settings}

\author{Moksh~Chadha%
\thanks{\authornote{Affiliation, e-mail and ORCID. Confirm affiliation policy before listing a company.}}}

\markboth{Preprint -- first draft}{Chadha: Edge VLMs for Person-Down Detection}

\maketitle

\begin{abstract}
A worker lying on the floor after a fall or collapse needs help within minutes, and camera-based alarms for this \pd{} event must run on cheap hardware close to the cameras. Small vision--language models (VLMs) can answer the question ``is a person down?'' without task-specific training, but it is unclear how small, how heavily quantized and how low-resolution they can be before they become unreliable. We study this trade-off with five open VLMs from 450\,M to 4\,B parameters, three quantization levels, three input resolutions, one to eight frames per decision and three prompts, and compare them with a pose-estimation baseline. We build a benchmark from four public fall datasets that separates true \pd{} events from look-alike postures such as kneeling, crouching, sitting and lying on furniture. We measure latency, peak memory and power on an NVIDIA Jetson Orin Nano and a CPU-only device and report the Pareto frontiers. \pending{headline result: smallest configuration within X F1 points of the best; false-alarm rate on hard negatives; streams per device at one decision per second}. A pilot study shows that the best prompt differs between models and that a detailed definition makes the smallest model answer ``yes'' to almost every frame. We release labels, code and per-run results.
\end{abstract}

\begin{IEEEkeywords}
Edge computing, fall detection, model quantization, occupational safety, vision--language models, zero-shot learning.
\end{IEEEkeywords}

\section{Introduction}
\label{sec:intro}

\IEEEPARstart{F}{alls} rank among the commonest causes of workplace death. In the United Kingdom, falls from height were the most frequent fatal accident to workers in 2025/26, about a quarter of all worker deaths \cite{hse2026fatal}; in the United States, falls, slips and trips caused 844 of the 5,070 fatal work injuries in 2024 \cite{bls2026cfoi}. An injured worker is often unable to summon help, and the outcome hinges on how swiftly someone notices. Wearable ``man-down'' alarms, which UK lone-working guidance recommends \cite{hse2020loneworkers}, mitigate this risk but depend on the worker wearing and charging a device. Fixed cameras already blanket many industrial floors, so detecting a person lying on the floor from video is a natural complement.

Two constraints govern such a system. It must run close to the cameras, since streaming every feed to a data centre consumes bandwidth, adds latency and moves footage of employees off site, which many sites will not tolerate; the detector must therefore fit on an inexpensive edge device serving several cameras. And it must separate a person who is down from postures that look alike in a single frame: a mechanic kneeling at a machine, a picker crouching at a low shelf, a worker sitting on the floor, someone lying under a vehicle. False alarms from such postures are why these systems end up ignored.

Classic fall detectors either learn a classifier from labelled videos \cite{mubashir2013survey,igual2013challenges} or apply hand-written rules to body keypoints \cite{cao2021openpose,maji2022yolopose}. Both excel on the datasets they were built for, yet both demand fresh labels or rules whenever the scene, camera angle or set of confusable postures shifts. Vision--language models (VLMs) offer an alternative: the task is stated in words, and the model answers without task-specific training \cite{radford2021clip,liu2023llava}. VLMs with under one billion parameters now run on edge hardware \cite{marafioti2025smolvlm,amini2025lfm2}, and integer quantization shrinks them further \cite{dettmers2022llmint8,frantar2023gptq,lin2024awq}. What is missing is evidence of where they cease to suffice. How much accuracy is forfeited in moving from 4\,B to 450\,M parameters, from 16-bit to 4-bit weights, or from full resolution to 256 pixels? Does a short clip beat a larger model? How many cameras can one device sustain?

We answer these questions for \pd{} detection in a zero-shot study: no model is fine-tuned, so the results describe what a practitioner obtains by downloading a model and writing a prompt. Our research questions are:
\begin{itemize}
  \item \textbf{RQ1} How do accuracy and false-alarm rate vary with model size, quantization, input resolution and number of frames per decision?
  \item \textbf{RQ2} What are the latency, peak memory and energy costs of each configuration on a consumer device, on its CPU alone and on its integrated GPU?
  \item \textbf{RQ3} Which configurations lie on the accuracy--latency and accuracy--memory Pareto frontiers?
  \item \textbf{RQ4} Can VLMs tell a person who is down from hard negatives such as kneeling, crouching, sitting on the floor and lying on furniture?
\end{itemize}

Our contributions are:
\begin{enumerate}
  \item A \pd{} benchmark of \BenchNKept{} frames and clips drawn from two public fall datasets \cite{kwolek2014urfall,auvinet2010mcfd}, with a written labelling policy, every frame checked by eye\iffullcheck\else\pending{after the author's full check}\fi{}, a blind agreement check, six hard-negative posture types, viewpoint annotations and sequence-disjoint development and test splits. We release labels and build scripts; the frames remain with their original distributors.
  \item A zero-shot evaluation of five open VLMs (450\,M--4\,B parameters) across quantization, resolution, temporal context and prompt wording, plus a keypoint-rule baseline that tests whether a VLM is needed at all.
  \item Latency, memory and energy measurements on the CPU and the integrated GPU of an Apple M4 Pro under a documented protocol, including two measurement pitfalls on unified-memory hardware that can distort reported memory by more than a factor of three.
  \item Accuracy--efficiency Pareto frontiers, a table of the camera streams each configuration sustains at one decision per second, and deployment recommendations for each device class.
\end{enumerate}

Section~\ref{sec:related} reviews related work. Section~\ref{sec:benchmark} describes the benchmark, Section~\ref{sec:methods} the models and baseline, and Section~\ref{sec:protocol} the evaluation protocol. Section~\ref{sec:results} reports results, Section~\ref{sec:discussion} discusses them with the study's limitations, and Section~\ref{sec:conclusion} concludes.

\section{Related Work}
\label{sec:related}

\subsection{Vision-based fall detection}
Fall detection has been studied for two decades, mostly for older people living alone \cite{mubashir2013survey,igual2013challenges}, and camera-based methods have been reviewed recently in \cite{gutierrez2021review,alam2022review}. Early methods used hand-crafted cues such as the aspect ratio of a person's silhouette, the orientation of a fitted ellipse and the height of the body above the floor; UR Fall's own baseline, for example, classifies depth-map features including the ratio of current to standing height \cite{kwolek2014urfall}. Deep networks trained on fall videos later replaced most of these features \cite{alam2022review}. A third line of work runs a pose estimator such as OpenPose \cite{cao2021openpose} or YOLO-Pose \cite{maji2022yolopose} and classifies the keypoints, either with learned models \cite{ramirez2021skeleton} or with rules on the speed of the hip centre, the angle of the body's centre line and the width-to-height ratio \cite{chen2020openposefall}. Our baseline follows the rule-based variant. Detector-based systems aimed at industrial sites also exist \cite{pereira2024industrialfall}.

Most of these methods are trained and tested on a few staged datasets \cite{kwolek2014urfall,martinez2019upfall,auvinet2010mcfd,charfi2013le2i}. OmniFall \cite{schneider2025omnifall} unifies eight of them under one 16-class taxonomy with cross-subject and cross-view splits, and adds synthetic and real-accident test data. Its results are directly relevant here: zero-shot multimodal models detected genuine falls about as well as fine-tuned models, but kept confusing a person who has \textit{fallen} with a person who is \textit{lying} down. Published results rarely break false alarms down by the posture that caused them; we do.

\subsection{Vision--language models and small VLMs}
Contrastive image--text pre-training \cite{radford2021clip} made zero-shot image classification practical, and connecting a vision encoder to a language model produced generative VLMs that answer free-form questions about images \cite{li2023blip2,liu2023llava}. Later work has pushed these models down to a few billion parameters or fewer for phones and embedded boards \cite{chu2023mobilevlm,beyer2024paligemma,marafioti2025smolvlm,bai2025qwen25vl,amini2025lfm2,qwen2026qwen35}, often with vision encoders that accept native aspect ratios and variable resolution \cite{tschannen2025siglip2}. The LFM2 family that we study was designed by searching architectures under on-device latency and memory constraints \cite{amini2025lfm2}. Model releases usually report general benchmarks and throughput on desktop or phone hardware; they say little about behaviour on one safety-critical yes/no question, or about cost on a particular edge board.

VLMs have known weaknesses on yes/no questions. They often answer ``yes'' when asked about objects that are not in the image \cite{li2023pope}, and this yes-bias is common enough to be studied as a failure mode in its own right \cite{guan2024hallusionbench,chan2026yesbias}. Their answers also change with superficial changes to the prompt \cite{zong2024permutations}. The bias does not always point the same way: in zero-shot surveillance anomaly detection, Yao et al.\ found that multimodal models favoured the ``normal'' class so strongly that recall collapsed, and that class-specific instructions moved the decision boundary \cite{yao2026surveillance}. For a \pd{} detector either direction is costly, as missed events or as false alarms, which is why we vary the prompt and report operating points chosen on data.

\subsection{VLMs for fall detection and workplace safety}
Few studies apply VLMs to fall detection. Junca et al.\ used a production multimodal model as the decision stage of a thermal-camera fall screener \cite{junca2026thermal}. They found that a single cautionary sentence in the prompt cost a large share of the detected falls, that identical requests at temperature zero still disagreed on 4\% of clips, and that removing the sentence led the model to call deliberate lie-downs falls, the same ambiguity our labelling policy addresses. Shahabian et al.\ compared zero-shot prompting and a few-shot probe on a VLM with skeleton-based and rule-based detectors for an assistive robot \cite{shahabian2026vlmfall}. OmniFall evaluated large zero-shot models on staged, synthetic and real falls \cite{schneider2025omnifall}.

In industrial safety, VLMs have been used to check personal protective equipment \cite{chen2024clip2safety}, detect safety violations in mining video \cite{wu2025monitorvlm}, recognise construction workers' actions \cite{bui2026construction} and inspect construction sites \cite{chen2026constructionsite10k}. SteelBench evaluates VLMs on real steel-plant CCTV, where workers appear small and are often occluded, and finds them well below human performance \cite{yarrabothula2026steelbench}. Closest to our setting, Adil et al.\ evaluated six VLMs under 4\,B parameters zero-shot for construction hazards and improved them by feeding object-detector output into the prompt \cite{adil2026svlm}.

These studies either use large or cloud-hosted models, or evaluate small models on accuracy alone. We found no study that measures how quantization, input resolution and temporal context trade accuracy against latency, memory and power for \pd{} detection on edge hardware, or that reports false alarms per look-alike posture. That is the gap this paper addresses.

\subsection{Quantization and on-device inference}
Post-training quantization reduces language-model weights to 8 or 4 bits with small accuracy loss \cite{dettmers2022llmint8,frantar2023gptq,lin2024awq}. llama.cpp \cite{gerganov2023llamacpp} implements block-wise integer formats in its GGUF file format and runs on CPUs and on Apple, NVIDIA and other GPUs, which makes it a common runtime for small models on edge devices. Laskaridis et al.\ benchmarked language models on phones and Jetson boards, including energy use and the accuracy cost of quantization \cite{laskaridis2024melt}, and Taherin et al.\ separated prefill from decode costs for llama.cpp and Transformers on several Jetson generations \cite{taherin2026hydra}. MLPerf defines standard latency and throughput scenarios for inference benchmarks \cite{reddi2020mlperf}. We follow this practice of reporting tail latency under a fixed scenario, and add two factors that matter for vision: the cost of encoding the image, which dominates latency when the answer is one word, and the effect of image resolution on both cost and accuracy.

\section{Benchmark}
\label{sec:benchmark}

\subsection{Source datasets}
No public dataset depicts falls in industrial settings under a research-permitting licence, so we build the benchmark from public indoor fall datasets (Table~\ref{tab:datasets}). Two complement each other. UR Fall \cite{kwolek2014urfall} offers per-frame posture labels, a camera parallel to the floor and, for falls, a ceiling-mounted second camera that affords the overhead view common in industrial installations. The Multiple Cameras Fall dataset (MCF) \cite{auvinet2010mcfd} films each of 24 scenarios, performed by one actor, with eight elevated wide-angle cameras, and includes sitting, lying on a sofa and housekeeping that supply hard negatives; we use the dense re-annotation published with OmniFall \cite{schneider2025omnifall}. We also intended to use UP-Fall \cite{martinez2019upfall} and Le2i \cite{charfi2013le2i}, but UP-Fall's public links returned a quota error and Le2i's repository was unavailable throughout; both are deferred. The two datasets used show staged events by actors in homes, offices or a laboratory (Section~\ref{sec:limitations}).

\begin{table*}[t]
\centering
\caption{Source datasets. Licence terms as stated on the official distribution pages on 20~October~2026.}
\label{tab:datasets}
\footnotesize
\setlength{\tabcolsep}{4pt}
\begin{tabular}{llllll}
\toprule
Dataset & Cameras & Resolution / fps & Viewpoints & Labels used & Licence \\
\midrule
UR Fall \cite{kwolek2014urfall} & 2 Kinect & $640{\times}480$ / 30 & side; ceiling (falls only) & per-frame lying / not lying (depth-based) & CC BY-NC-SA 4.0 \\
MCF \cite{auvinet2010mcfd} & 8 IP, wide-angle & $720{\times}480$ / 30$^\ast$ & elevated, oblique & OmniFall segments \cite{schneider2025omnifall} & none stated \\
\bottomrule
\end{tabular}

\smallskip
\raggedright\footnotesize $^\ast$The AVI headers state 120\,fps; frame counts match OmniFall's timeline at 30\,fps, which we use.
\end{table*}

\subsection{Labelling policy}
Each sample receives one of three labels. A sample is \pd{} when a person's torso lies on the floor or ground, however it got there; this covers falls and deliberate lying on the floor. UR Fall counts deliberate lying as an activity of daily living, hence a negative; a single frame cannot distinguish the two, and for a safety alarm both warrant a check, so we label them \pd{} with the subtype \textit{lying\_floor\_deliberate} and report them separately. Junca et al.\ met the same ambiguity: a prompt that caught more falls also flagged deliberate lie-downs \cite{junca2026thermal}.

A sample is a \textit{hard negative} when a person adopts a low or horizontal posture with the torso off the floor: crouching, kneeling (including on hands and knees), bending over, sitting (chair, sofa, bed or floor), or lying on a sofa or bed. A sample is \textit{normal} when every visible person stands or walks. Labels describe the image, not the actor's activity, so a frame with one person standing and another kneeling beside a mattress is a kneeling hard negative. We exclude transitional frames (mid-fall, getting down or up), frames with no visible person, and frames whose posture is equivocal, such as a fallen person propped up or a torso draped over a table.

\subsection{Construction}
Candidates are proposed automatically and each is reviewed visually. For UR Fall they derive from the official per-frame depth labels and the ratio of current to standing height; the overhead camera, unlabelled, inherits the label of the side-camera frame with the same index (both streams have equal frame counts). For MCF, candidates are sampled inside OmniFall segments labelled \textit{fallen}, \textit{sitting}, \textit{lying}, \textit{standing}, \textit{walk} and \textit{other}, at least 0.25\,s from segment boundaries. Per source sequence we take at most three frames per label (MCF: two per scenario, camera and label, one for walking), at least 1\,s apart, then subsample MCF to a fixed quota per label.

An AI assistant (Claude, Anthropic) conducted the visual review under the author's direction, inspecting contact sheets of every candidate and recording an accept, reject or relabel decision, with a reason, for each frame, per the policy above. Because a vision--language model helped label a benchmark for vision--language models, we validated the labels twice. First, a blind check: \HumanCheckText{} Second, \iffullcheck the author went through every frame with its label displayed, confirming \FullConfirmed{}, changing \FullChanged{} and removing \FullRemoved{}. A visible label invites anchoring, so the blind check is the unbiased measure; the frames the author labelled differently in the blind check were all confirmed here, so the full pass is a screen for gross errors, not independent relabelling.\else\pending{every frame checked by eye by the author: confirmed / changed / removed counts}\fi{} Of \BenchNCand{} candidates, the review kept \BenchNKept{}, rejected \BenchNRejected{}, and relabelled \BenchNRelabelled{} of those kept. These corrections matter to anyone reusing the datasets: several UR Fall frames labelled ``lying on the ground'' show a person on hands and knees, ``upright'' frames include people bending over a sofa, and some MCF segments labelled \textit{sitting} show an empty view from cameras that cannot see the chair.

Every sample also carries a clip: the labelled frame plus the seven preceding frames at 0.25\,s spacing (1.75\,s in all), so decisions on 1, 4 or 8 frames end at the same labelled frame. Each sample records its viewpoint: \textit{side} (UR Fall camera~0), \textit{overhead} (UR Fall camera~1) or \textit{elevated} (MCF).

\subsection{Splits}
Frames from one sequence are strongly correlated, and MCF films one moment from eight cameras. We therefore split by sequence: an MCF scenario with all its cameras, or a UR Fall sequence with both, goes wholly to development or wholly to test. About a quarter of the sequences of each kind form the development set, which serves solely to choose thresholds. The test set has \BenchNTest{} samples from \BenchNSeqTest{} sequences (\BenchTestDown{} \pd{}, \BenchTestHard{} hard negative, \BenchTestNormal{} normal), the development set \BenchNDev{} samples from \BenchNSeqDev{} sequences (Table~\ref{tab:benchmark}). Because UR Fall's overhead camera recorded only falls, overhead frames are almost uniformly \pd{}: for that viewpoint we report recall but not a false-alarm rate.

\begin{table}[t]
\centering
\caption{Benchmark composition after visual review.}
\label{tab:benchmark}
\footnotesize
\setlength{\tabcolsep}{2.4pt}
\begin{tabular}{llrrrr}
\toprule
Label & Subtype & UR Fall & MCF & Test & Dev \\
\midrule
person down & fall & 82 & 136 & 160 & 58 \\
 & lying on floor, deliberate & 15 & 0 & 10 & 5 \\
\midrule
hard negative & crouching & 11 & 20 & 22 & 9 \\
 & kneeling (incl.\ hands and knees) & 10 & 11 & 16 & 5 \\
 & bending & 14 & 29 & 25 & 18 \\
 & sitting on a chair & 8 & 53 & 47 & 14 \\
 & sitting on a sofa or bed & 9 & 41 & 19 & 31 \\
 & sitting on the floor & 1 & 0 & 1 & 0 \\
 & lying on a sofa or bed & 18 & 38 & 41 & 15 \\
\midrule
normal & standing & 62 & 33 & 65 & 30 \\
 & walking & 0 & 120 & 88 & 32 \\
\midrule
Total & & 230 & 481 & 494 & 217 \\
\bottomrule
\end{tabular}

\end{table}

\subsection{Reproducibility}
UR Fall is licensed for non-commercial use with share-alike terms and MCF states no licence, so we do not redistribute frames. We release the label files (dataset, sequence, camera, frame index, label, subtype, viewpoint, review note) and scripts that download the source data from the official pages and rebuild every sample. The UR Fall script fetches only the needed frames through HTTP range requests, keeping the download below 1\,GB.

\section{Models and Methods}
\label{sec:methods}

\subsection{Task formulation}
Given one image, or $n \in \{4, 8\}$ consecutive frames ending at the labelled frame, the system decides whether a person is down. A VLM receives the image(s) and a prompt and emits a short answer, mapped to \textit{yes}, \textit{no} or \textit{unparsed} (Section~\ref{sec:parsing}). Because a hard yes/no pins the operating point wherever the model happens to prefer, we also read its probability of ``yes'': the probability mass on ``yes'' tokens at the first generated position, renormalised over ``yes'' and ``no''. This score, $p_\text{yes}$, permits threshold-free metrics and lets us choose operating points (Section~\ref{sec:thresholds}).

\subsection{Models}
Table~\ref{tab:models} lists the models. The LFM2.5-VL family \cite{amini2025lfm2,liquidai2026lfm25vl450m,liquidai2026lfm25vl3b} spans three sizes built alike: a SigLIP2 NaFlex vision encoder \cite{tschannen2025siglip2} (86\,M parameters in the smallest model, 400\,M in the others) feeds a hybrid language model of 350\,M, 1.2\,B or 2.6\,B parameters. Varying size within one family holds the training recipe and preprocessing roughly constant. Qwen3.5 at 2\,B and 4\,B \cite{qwen2026qwen35} serves as a reference from another family. We chose models that run in llama.cpp \cite{gerganov2023llamacpp} on our device and carry research-permitting licences. Liquid AI publishes official GGUF exports; Qwen does not, so for Qwen3.5 we use Unsloth's community conversions. Qwen3.5 emits a reasoning trace by default; we disable it so that every model answers directly. Appendix~\ref{app:revisions} lists exact model revisions. We verified that our llama.cpp build accepts image input for both families.

\begin{table*}[t]
\centering
\caption{Models evaluated. Disk size is Q8\_0 language-model weights plus the F16 vision projector (GGUF). $^\dagger$Community GGUF conversion (Unsloth); Qwen publishes no official GGUF.}
\label{tab:models}
\footnotesize
\begin{tabular}{lrlllr}
\toprule
Model & Params & Language backbone & Vision encoder & Licence & Disk (MiB) \\
\midrule
LFM2.5-VL-450M & 0.45\,B & LFM2.5-350M (hybrid) & SigLIP2 NaFlex, 86\,M & LFM Open License 1.0 & 362 + 180 \\
LFM2.5-VL-1.6B & 1.6\,B & LFM2.5-1.2B (hybrid) & SigLIP2 NaFlex, 400\,M & LFM Open License 1.0 & 1189 + 814 \\
LFM2.5-VL-3B & 3\,B & LFM2.5-2.6B (hybrid) & SigLIP2 NaFlex, 400\,M & LFM Open License 1.0 & 2742 + 814 \\
Qwen3.5-2B$^\dagger$ & 2.27\,B & hybrid linear/full attention, 24 layers & ViT, 24 layers & Apache 2.0 & 1919 + 637 \\
Qwen3.5-4B$^\dagger$ & 4.66\,B & hybrid linear/full attention, 32 layers & ViT, 24 layers & Apache 2.0 & 4275 + 641 \\
\bottomrule
\end{tabular}
\end{table*}

\subsection{Prompts}
We test three prompts (full text in Appendix~\ref{app:prompts}). Prompt~\promptA{} is a plain question: ``Is any person fallen or lying on the floor? Answer yes or no.'' Prompt~\promptB{} adds a definition of \pd{} and enumerates the postures that do not count, from kneeling and crouching to lying on furniture or under equipment. Prompt~\promptBtwo{} gives only a one-sentence criterion: ``A person is `down' only if their torso is lying on the floor or ground. Is any person down? Answer yes or no.'' We added \promptBtwo{} after the pilot study (Section~\ref{sec:pilot}) showed that listing look-alike postures in \promptB{} increased false alarms; its criterion mirrors the pose baseline's rule. For multi-frame inputs, a fixed prefix states that the images are consecutive frames in time order and asks about the last.

\subsection{Runtimes and quantization}
All main-grid measurements use llama.cpp \cite{gerganov2023llamacpp} with GGUF weights at three precisions: 16-bit (F16, or BF16 for Qwen), Q8\_0 (8-bit) and Q4\_K\_M (mostly 4-bit, some tensors at 6-bit). The vision projector stays 16-bit throughout, so only the language model's precision varies. We run \texttt{llama-server} with one slot, a 4096-token context and greedy decoding, via its OpenAI-compatible API. In the pilot we also ran LFM2.5-VL-450M in bf16 with Hugging Face Transformers on the same device as an unquantized cross-check. We do not evaluate ONNX Runtime.

\subsection{Input resolution}
\label{sec:resolution}
We resize each frame so that its longest side measures the native size, 512 or 256 pixels, with Lanczos filtering, never upscaling. Frames are sent as lossless PNG, so resolution is the only image variable. The model's own preprocessor resizes and tiles the image again before encoding, so the number of image tokens depends on both our resizing and the model; we report it for every configuration.

\subsection{Output parsing}
\label{sec:parsing}
We cap generation at 16 tokens. An answer is \textit{strictly} parsed if it begins with ``yes'' or ``no'' (any case, optional punctuation) and \textit{leniently} parsed if exactly one of the two words occurs in its first sentence; anything else is \textit{unparsed} and counted as an error. The gap between strict and lenient parsing gauges how faithfully each model follows the format instruction.

\subsection{Decision thresholds}
\label{sec:thresholds}
We report two operating points. The \textit{default} uses the generated answer, equivalent to thresholding $p_\text{yes}$ at roughly 0.5. The \textit{calibrated} point thresholds $p_\text{yes}$ at the value attaining 95\% recall on the development set, because a safety alarm should seldom miss a real event. The threshold consumes labelled data but trains no parameters; we report both so readers can see how much of each model's error originates in its default bias.

\subsection{Pose-based baseline}
\label{sec:baseline}
To test whether a VLM is needed at all, we compare against a keypoint rule. YOLO11 pose models \cite{maji2022yolopose,jocher2024yolo11} (nano and small) estimate 17 body keypoints and a bounding box per person. A person is flagged as down when (i) the torso axis, from hip midpoint to shoulder midpoint, lies within $\theta$ degrees of horizontal and (ii) the bounding box is at least $r$ times as wide as tall; for clips, both must hold in at least a fraction $q$ of the frames. Similar rules have been applied to fall detection before \cite{chen2020openposefall}. We tune $\theta \in \{15^\circ, 20^\circ, \dots, 60^\circ\}$, $r \in \{0.6, 0.8, \dots, 2.0\}$ and, for clips, $q$ on the development set only, with the VLMs' calibrated criterion: the best development F1 among settings with at least 95\% development recall, or the best F1 overall if none reaches that recall. The rule needs no language model and costs a fraction of a VLM's compute, so the question is whether the VLMs' greater robustness to hard negatives justifies their latency.

\section{Evaluation Protocol}
\label{sec:protocol}

\subsection{Accuracy metrics}
We treat \pd{} as the positive class and report precision, recall and F1 at both operating points, the area under the ROC curve (AUROC) and the area under the precision--recall curve computed from $p_\text{yes}$, and the false-alarm rate (FAR), the fraction of negatives answered ``yes''. FAR is reported separately for normal frames, for all hard negatives and for each hard-negative subtype, because a single pooled number hides which postures cause alarms. We also report accuracy by camera viewpoint. Confidence intervals are 95\% percentile intervals from a cluster bootstrap that resamples whole source sequences \cite{efron1993bootstrap}, since frames from one sequence are not independent. We use 10,000 bootstrap replicates. When a proportion is 0 or 1 the percentile bootstrap collapses to a zero-width interval, so in that case we report the Wilson score interval instead.

\subsection{Efficiency metrics}
For each configuration we measure:
\begin{itemize}
  \item \textit{End-to-end latency}: wall-clock time from sending the request to receiving the last token, measured by the client. This includes image encoding by the client, transfer to the server, vision encoding, prompt processing and decoding.
  \item \textit{Time to first token} (TTFT), measured by the client from a streamed response, and the server-side prefill and decode times reported by llama.cpp.
  \item \textit{Decode throughput} in tokens per second. With one-word answers decoding is a small share of latency, but it matters for the bounding-box variant.
  \item \textit{Peak memory} (see below), \textit{model size on disk} (language model plus vision projector), and \textit{power}.
  \item \textit{Streams per device}: the number of cameras one device can serve at one decision per second per camera, computed as $\lfloor 1000 / t_{95} \rfloor$ where $t_{95}$ is the 95th-percentile end-to-end latency in milliseconds, for sequential processing. We do not batch requests from different cameras; batching can raise throughput at the cost of latency, and we leave it to future work.
\end{itemize}

\textit{Measuring memory.} Peak memory is harder to measure than it looks, and in our pilot two common metrics disagreed by more than a factor of three. Resident set size (RSS) counts pages of the memory-mapped weight file once they are touched. The operating system can evict these pages and reload them from disk, so RSS can overstate the memory a process truly needs. On Apple Silicon the ``physical footprint'' metric excludes those file-backed pages but also excludes GPU buffers that wrap them. For the 3\,B model on the Metal backend it reported \PilotLargeGpuFootprintGiB\,GiB, less than half the size of the weights on disk. On the CPU backend llama.cpp repacks quantized weights into a CPU-friendly layout and keeps the original mapping, which roughly doubles RSS (\PilotLargeCpuRssGiB\,GiB for the 3\,B model with 3.5\,GiB of weights). Table~\ref{tab:memory} separates these effects on the pilot machine. Disabling memory mapping (\texttt{--no-mmap}) copies the weights into private memory, after which RSS and footprint agree to within 0.06\,GiB in every configuration we tried: \AblLargeGpuNoMmapRssGiB\ and \AblLargeGpuNoMmapFpGiB\,GiB for the 3\,B model on the GPU. Disabling repacking (\texttt{--no-repack}) removes the second copy on the CPU and cuts the 3\,B model's RSS from \AblLargeCpuDefaultRssGiB\ to \AblLargeCpuNoRepackRssGiB\,GiB, at the price of slower CPU inference. None of the flags changed a single answer. We also tried a system-level measure, the rise in total-minus-available memory over an idle baseline, but on macOS it was not reproducible: with memory compression and background processes it sometimes reported less than the process's own private memory. We therefore run all efficiency measurements with \texttt{--no-mmap} and report peak RSS, which then equals the memory the model actually occupies; on the Jetson we additionally report total used RAM from \texttt{tegrastats}, since CPU and GPU share memory there. Repacking stays enabled, as it would in a deployment, and its memory cost is visible in the CPU numbers.

\begin{table*}[t]
\centering
\caption{Memory accounting on the pilot machine (\PilotCpu, \PilotRamGB\,GB, macOS), Q8\_0, prompt~\promptA, 30 requests per row. ``Same answers'' compares with the default flags. Latencies are relative only: other applications were running.}
\label{tab:memory}
\footnotesize
\setlength{\tabcolsep}{5pt}
\begin{tabular}{lllrrrrc}
\toprule
Model & Backend & Flags & Load & Median & RSS & Footpr. & Same \\
 & & & (s) & (ms) & (GiB) & (GiB) & answers \\
\midrule
LFM2.5-VL-450M & CPU & \texttt{default} & 0.4 & 1020 & 1.09 & 0.71 & 30/30 \\
 & CPU & \texttt{-{}-no-mmap} & 0.4 & 1125 & 0.87 & 0.84 & 30/30 \\
 & CPU & \texttt{-{}-no-repack} & 0.4 & 1295 & 0.77 & 0.38 & 30/30 \\
 & CPU & \texttt{-{}-no-mmap -{}-no-repack} & 0.4 & 1365 & 0.76 & 0.73 & 30/30 \\
\midrule
LFM2.5-VL-450M & Metal & \texttt{default} & 0.4 & 119 & 0.70 & 0.40 & 30/30 \\
 & Metal & \texttt{-{}-no-mmap} & 0.4 & 120 & 0.84 & 0.89 & 30/30 \\
\midrule
LFM2.5-VL-3B & CPU & \texttt{default} & 2.7 & 5071 & 6.42 & 3.72 & 30/30 \\
 & CPU & \texttt{-{}-no-mmap} & 1.7 & 5471 & 4.29 & 4.24 & 30/30 \\
 & CPU & \texttt{-{}-no-repack} & 0.8 & 6270 & 3.92 & 1.20 & 30/30 \\
 & CPU & \texttt{-{}-no-mmap -{}-no-repack} & 1.0 & 6161 & 3.92 & 3.86 & 30/30 \\
\midrule
LFM2.5-VL-3B & Metal & \texttt{default} & 0.8 & 576 & 3.72 & 1.09 & 30/30 \\
 & Metal & \texttt{-{}-no-mmap} & 1.2 & 560 & 4.25 & 4.27 & 30/30 \\
\bottomrule
\end{tabular}

\end{table*}

\textit{Measuring power.} On the Jetson we sample total board power with \texttt{tegrastats} at 100\,ms intervals and report mean power during inference and energy per decision. \pending{CPU device: external USB/wall power meter, model and sampling rate}.

\subsection{Hardware and software}
Table~\ref{tab:hardware} lists the devices. The full grid runs on the Jetson Orin Nano; the CPU-only device runs only the configurations that are Pareto-optimal on the Jetson. Accuracy runs that do not depend on the device run on a cloud GPU; because decoding is greedy and deterministic, a configuration gives the same answers on every backend we tested (Section~\ref{sec:pilot}), so accuracy measured once applies to all devices.

\begin{table}[t]
\centering
\caption{Hardware and software. Jetson specifications are for the Orin Nano Super Developer Kit \cite{nvidia_orin_nano_super}. \pending{software rows; CPU-only device}}
\label{tab:hardware}
\setlength{\tabcolsep}{3pt}
\begin{tabular}{lll}
\toprule
 & Jetson Orin Nano & CPU-only device \\
\midrule
CPU & 6$\times$ Cortex-A78AE, 1.7\,GHz & \pending{} \\
GPU & Ampere, 1024 CUDA cores & none (CPU only) \\
Memory & 8\,GB LPDDR5, 102\,GB/s & \pending{} \\
OS / JetPack & \pending{} & \pending{} \\
Power mode & \pending{nvpmodel} & \pending{} \\
Threads & \pending{} & \pending{} \\
llama.cpp build & \pending{} & \pending{} \\
\bottomrule
\end{tabular}
\end{table}

\subsection{Measurement procedure}
\label{sec:procedure}
Each configuration runs in a fresh server process. We discard three warm-up requests and then time at least 50 requests, reporting the median and the 95th percentile. Prompt caching is disabled, so every request processes its image and prompt from scratch, as a new camera frame would. Requests are sent one at a time from the same machine. The device runs nothing else, is connected to mains power, and is kept from sleeping. We added the last condition after a pilot run in which macOS put the machine into idle sleep in the middle of a request; that request was logged as taking 1015\,s instead of about 5\,s. We record the operating system, runtime build, thread count, power mode and 1-minute load average with every run.

\section{Results}
\label{sec:results}

\subsection{Pilot study}
\label{sec:pilot}
Before running the full grid we tested the pipeline on 30 frames from UR Fall: 10 \pd{} (five from the side camera, five from the ceiling camera), 10 hard negatives (crouching, kneeling, bending, sitting on a chair or bed, lying on a bed or sofa) and 10 normal frames, one frame per sequence and every frame checked by eye. We ran LFM2.5-VL-450M, LFM2.5-VL-3B and Qwen3.5-2B at Q8\_0 with all three prompts. The pilot is too small for conclusions about the models: the 3\,B model's hard-negative false-alarm rate with prompt~\promptA{}, 2/10, has a 95\% bootstrap interval of \PilotLargeAFaHardCI, and a perfect 10/10 recall has a 95\% Wilson interval of \PilotRecallWilsonCI. We report it because it shaped the protocol and the prompt set.

\begin{table}[t]
\centering
\caption{Pilot study on 30 UR Fall frames (10 per class), default operating point, Q8\_0. FA = false alarms. AUROC is computed from $p_\text{yes}$. Too small for model comparisons; reported because it shaped the protocol.}
\label{tab:pilot}
\footnotesize
\setlength{\tabcolsep}{3.5pt}
\begin{tabular}{llccccc}
\toprule
Model (Q8\_0) & Prompt & Recall & FA$_\text{hard}$ & FA$_\text{normal}$ & F1 & AUROC \\
\midrule
LFM2.5-VL-450M & A & 10/10 & 5/10 & 0/10 & 0.80 & 1.00 \\
 & B & 10/10 & 9/10 & 10/10 & 0.51 & 0.89 \\
 & B2 & 10/10 & 9/10 & 9/10 & 0.53 & 0.94 \\
\midrule
LFM2.5-VL-3B & A & 10/10 & 2/10 & 0/10 & 0.91 & 1.00 \\
 & B & 10/10 & 6/10 & 0/10 & 0.77 & 1.00 \\
 & B2 & 10/10 & 3/10 & 0/10 & 0.87 & 1.00 \\
\midrule
Qwen3.5-2B & A & 10/10 & 2/10 & 0/10 & 0.91 & 1.00 \\
 & B & 10/10 & 4/10 & 0/10 & 0.83 & 1.00 \\
 & B2 & 10/10 & 0/10 & 0/10 & 1.00 & 1.00 \\
\bottomrule
\end{tabular}

\end{table}

Four observations carried over into the design. First, output format was not a problem: all \PilotNAnswers{} answers across models, prompts and backends began with ``yes'' or ``no'' (\PilotNLenient{} needed lenient parsing, \PilotNUnparsed{} were unparsed), so we did not need constrained decoding. Second, prompt wording changed the false-alarm rate more than model size did (Table~\ref{tab:pilot}). The detailed definition in prompt~\promptB{} raised the 3\,B model's hard-negative false alarms from \PilotLargeAFaHard{} to \PilotLargeBFaHard{} and made the 450\,M model answer ``yes'' to \PilotSmallBFaNormal{} frames of people simply standing or walking. Listing the excluded postures seems to prime the model toward them. This led us to add prompt~\promptBtwo{}, which defines \pd{} by one physical criterion and names no excluded posture. The effect of \promptBtwo{} depended on the model. It removed every hard-negative false alarm for Qwen3.5-2B (\PilotQwenBtwoFaHard{}, against \PilotQwenAFaHard{} with \promptA), brought the 3\,B LFM model back to \PilotLargeBtwoFaHard{}, and did not help the 450\,M model (\PilotSmallBtwoFaNormal{} normal frames flagged), whose $p_\text{yes}$ on standing people stayed just above 0.5. No single prompt was best for all three models. Third, with prompt~\promptA{} all three models ranked every \pd{} frame above every negative (AUROC \PilotLargeAAuroc{}), so their errors came from where the default answer places the threshold rather than from an inability to separate the classes. This is why we report a calibrated operating point. Fourth, answers were identical across backends: \PilotCpuGpuAgree{} answers matched between CPU and Metal, and \PilotQeightBfAgree{} between Q8\_0 GGUF and bf16 Transformers for the 450\,M model, the one difference being a frame with $p_\text{yes}$ close to 0.5.

The labelling check also mattered. UR Fall's depth-based annotation marks a frame as ``lying on the ground'' from the person's height above the floor. Viewing the candidates showed that the ``lying'' frames in its activity-of-daily-living sequences mix lying on a bed or sofa with deliberate lying on the floor. Under our policy the first are hard negatives and the second are \pd{}. All three models flagged all nine deliberate floor-lying frames under every prompt (\PilotFloorHits{} answers). These frames were filmed in rooms other than the one where all UR Fall falls take place, and with prompt~\promptA{} the same models rejected the standing people filmed in those rooms (\PilotSameRoomNormalRejectA{} answers), which is early evidence that the models respond to posture rather than recognising the fall room.

On an \PilotCpu{} laptop, used here only to validate the measurement pipeline, a decision took a median of \PilotSmallCpuMedianMs\,ms (LFM 450\,M), \PilotLargeCpuMedianMs\,ms (LFM 3\,B) and \PilotQwenCpuMedianMs\,ms (Qwen 2\,B) on 10 CPU threads, and \PilotSmallGpuMedianMs, \PilotLargeGpuMedianMs{} and \PilotQwenGpuMedianMs\,ms on the integrated GPU (Table~\ref{tab:pilot_eff}). Qwen3.5-2B encodes the same frame into more tokens than the LFM models, which offsets its smaller language model. On the CPU, prefill (vision encoding and prompt processing) accounted for \PilotSmallCpuPrefillPct--\PilotLargeCpuPrefillPct\% of the time while decoding the one-word answer took a median of \PilotSmallCpuDecodeMs\,ms (450\,M) and \PilotLargeCpuDecodeMs\,ms (3\,B), which suggests that input resolution will be the main lever on latency.

\begin{table}[t]
\centering
\caption{Pilot efficiency on an \PilotCpu{} (\PilotRamGB\,GB, macOS \PilotOs, llama.cpp b\PilotLlamaBuild), prompt~\promptA, 30 timed requests after 3 warm-ups. Validates the pipeline only; edge results are in Table~\ref{tab:efficiency}.}
\label{tab:pilot_eff}
\footnotesize
\setlength{\tabcolsep}{3pt}
\begin{tabular}{llrrrr}
\toprule
Model (Q8\_0) & Backend & Median & p95 & Peak RSS & Streams \\
 & & (ms) & (ms) & (GiB) & @1\,Hz \\
\midrule
LFM2.5-VL-450M & CPU, 10 thr. & 916 & 972 & 1.08 & 1 \\
LFM2.5-VL-3B & CPU, 10 thr. & 4270 & 4464 & 6.42 & 0 \\
Qwen3.5-2B & CPU, 10 thr. & 4967 & 5132 & 4.75 & 0 \\
\midrule
LFM2.5-VL-450M & Metal GPU & 112 & 116 & 0.70 & 8 \\
LFM2.5-VL-3B & Metal GPU & 530 & 537 & 3.73 & 1 \\
Qwen3.5-2B & Metal GPU & 523 & 536 & 2.82 & 1 \\
\bottomrule
\end{tabular}

\end{table}

\subsection{RQ1: accuracy versus model size, quantization, resolution and frames}
\pending{Main accuracy table: one row per configuration on the test split, with F1, recall, FAR (hard negatives, normal) at default and calibrated operating points, AUROC, 95\% CIs. Then marginal plots: F1 versus parameters (one line per quantization), F1 versus resolution, F1 versus frames per decision.}

\pending{Text: size effect within LFM2.5-VL; quantization effect (expect Q8\_0 $\approx$ F16; Q4\_K\_M cost per size); resolution effect and its interaction with model size; whether 4/8 frames help and whether they help more on hard negatives.}

\subsection{RQ2: latency, memory and power on edge devices}
\begin{table*}[t]
\centering
\caption{Efficiency on the edge devices. \pending{generate from results}}
\label{tab:efficiency}
\begin{tabular}{lllrrrrrrr}
\toprule
Device & Model & Quant & Res. & Median (ms) & p95 (ms) & TTFT (ms) & Peak RSS (GiB) & Power (W) & Streams \\
\midrule
\multicolumn{10}{c}{\pending{rows}} \\
\bottomrule
\end{tabular}
\end{table*}

\pending{Text: latency scaling with image tokens; memory floor per model; energy per decision; streams per device at 1\,Hz.}

\subsection{RQ3: Pareto frontiers}
\begin{figure}[t]
\centering
\fbox{\parbox[c][4.5cm][c]{0.9\columnwidth}{\centering\pending{Figure: F1 (calibrated) versus p95 latency on Jetson Orin Nano; Pareto-optimal points highlighted; pose baseline marked.}}}
\caption{Accuracy--latency trade-off on the Jetson Orin Nano. \pending{}}
\label{fig:pareto_latency}
\end{figure}

\begin{figure}[t]
\centering
\fbox{\parbox[c][4.5cm][c]{0.9\columnwidth}{\centering\pending{Figure: F1 (calibrated) versus peak memory.}}}
\caption{Accuracy--memory trade-off. \pending{}}
\label{fig:pareto_memory}
\end{figure}

\pending{Text: which configurations are Pareto-optimal on each device; where the pose baseline sits.}

\subsection{RQ4: hard negatives}
\begin{table}[t]
\centering
\caption{False-alarm rate (\%) per hard-negative subtype for the Pareto-optimal configurations and the pose baseline. \pending{}}
\label{tab:hardneg}
\begin{tabular}{lccccc}
\toprule
Config & Crouch & Kneel & Bend & Sit & Lie (furn.) \\
\midrule
\multicolumn{6}{c}{\pending{rows}} \\
\bottomrule
\end{tabular}
\end{table}

\pending{Text: which subtypes cause most alarms (pilot suggests kneeling / hands-and-knees and crouching); effect of prompt and of temporal context on each subtype; accuracy by viewpoint (overhead vs. oblique).}

\subsection{Bounding-box output}
\pending{For the best configuration only: prompt for a box around the person who is down; report how often the box overlaps the annotated person (IoU $\geq 0.5$) and the latency cost of the longer output.}

\subsection{Failure cases}
\begin{figure*}[t]
\centering
\fbox{\parbox[c][4cm][c]{0.95\textwidth}{\centering\pending{Figure: grid of missed falls and hard-negative false alarms for the recommended configurations. Faces blurred. Caption must state the source dataset licence (UR Fall and Le2i frames: CC BY-NC-SA).}}}
\caption{Typical failures. \pending{}}
\label{fig:failures}
\end{figure*}

\section{Discussion}
\label{sec:discussion}

\subsection{Deployment recommendations}
The results support a concise set of recommendations. In every case, use the plain question, one frame per decision, and a threshold on $p_\text{yes}$ chosen on labelled footage from the site rather than the model's own yes/no answer.

\textit{Device with an integrated GPU or accelerator.} LFM2.5-VL-1.6B with Q4\_K\_M weights and 256\,px input offers the best balance we measured: F1 \CLfmMQfourTwofiftysixOneFone, recall \CLfmMQfourTwofiftysixOneRecall\%, hard-negative false alarms \CLfmMQfourTwofiftysixOneFarHard\%, p95 latency \CLfmMQfourTwofiftysixOnePtf\,ms and \CLfmMQfourTwofiftysixOneRss\,GiB, sufficient for \CLfmMQfourTwofiftysixOneStreams{} cameras at one decision per second each, at \PLfmMQfourTwofiftysixOneGpuJ\,J per decision. Where accuracy outweighs camera count, LFM2.5-VL-3B with Q8\_0 weights at 512\,px reaches F1 \CLfmLQeightFivetwelveOneFone{} with \CLfmLQeightFivetwelveOneFarHard\% hard-negative false alarms, at \CLfmLQeightFivetwelveOnePtf\,ms and \CLfmLQeightFivetwelveOneRss\,GiB.

\textit{CPU only.} Only LFM2.5-VL-450M sustains one decision per second for one camera at useful accuracy: with Q8\_0 weights at 512\,px it reaches F1 \CLfmSQeightFivetwelveOneFone{} (recall \CLfmSQeightFivetwelveOneRecall\%, hard-negative false alarms \CLfmSQeightFivetwelveOneFarHard\%) at \CLfmSQeightFivetwelveOneCpuPtf\,ms, \CLfmSQeightFivetwelveOneRss\,GiB and \PLfmSQeightFivetwelveOneCpuJ\,J per decision. If a slower check rate is tolerable, LFM2.5-VL-1.6B at 256\,px gives F1 \CLfmMQfourTwofiftysixOneFone{} at one decision every \CLfmMQfourTwofiftysixOneCpuPtf\,ms per camera. Our CPU is a fast laptop core; slower embedded CPUs will need the 450\,M model or a lower decision rate. A GPU or accelerator, where available, should be used: the same configuration consumed \PCpuGpuRatioMin--\PCpuGpuRatioMax{} times more energy per decision on the CPU.

\textit{What to avoid.} Long, definition-style prompts; 256\,px input with any model but LFM2.5-VL-1.6B; multi-frame inputs, which multiply latency without improving accuracy; and 16-bit weights, which occupy about twice the memory of Q4\_K\_M with no consistent accuracy gain.

\subsection{Prompting small models}
A prompt that defines the task meticulously and enumerates the cases to exclude reads better to a human, yet in our grid it drove the 450\,M model to answer ``yes'' to nearly every frame and raised false alarms for every LFM model (Table~\ref{tab:prompts}). VLMs are known to answer ``yes'' too readily \cite{li2023pope,chan2026yesbias}, and in thermal fall screening a single cautionary sentence shifted sensitivity substantially \cite{junca2026thermal}. Our results suggest that naming the excluded postures amplifies the yes-bias of small models instead of correcting it. A one-sentence criterion (\promptBtwo) looked best for Qwen3.5-2B in the 30-frame pilot, but on the full benchmark it rendered that model conservative, and it missed about half of the falls. Across all five models the plain question was best or tied. The pilot is a salutary reminder of how misleading 30 samples can be.

\subsection{Thresholds instead of answers}
The default answer is the weakest element of these models. Their scores rank falls above look-alike postures well (AUROC \RLfmSAuroc--\RLfmLAuroc), but the cut-off implied by the yes/no answer lies elsewhere for each model: LFM2.5-VL-450M answered ``yes'' for \RLfmSDefFarHard\% of hard negatives, Qwen3.5-4B for \RQwenLDefFarHard\%. Choosing the threshold on development data reduced false alarms for every model, most for the smallest. Token probabilities are not necessarily calibrated \cite{guo2017calibration}, so the threshold must be learned from labelled data, and it does not transfer perfectly: two models reached 90--91\% recall on the test set with thresholds set for 95\% on development. Reading token probabilities is free in llama.cpp and most runtimes, so deployments should log $p_\text{yes}$ and set, and periodically re-check, a site-specific threshold.

\subsection{Is a VLM needed?}
For this task, yes. The pose-and-rule baseline reached F1 \CPoseNPoseNatOneFone--\CPoseNPoseNatFourFone, against at least \RLfmSCalFone{} for every VLM, because the pose detector frequently fails to find a person lying on the floor at all. Its speed (\RPoseNMs\,ms per frame) is attractive, but an inexpensive first stage that overlooks about half of the fallen people cannot gate a slower model. A pose model trained on fallen people, or a detector-guided prompt as in \cite{adil2026svlm}, might alter this verdict; we leave both to future work.

\subsection{Measuring efficiency}
Two details of our procedure changed results by large factors. The first is memory accounting on unified-memory hardware: depending on the metric and on two runtime flags, the same 3\,B model ``used'' between \AblLargeGpuDefaultFpGiB\ and \AblLargeCpuDefaultRssGiB\,GiB (Table~\ref{tab:memory}). Papers that report one peak-memory figure without naming the metric and runtime settings resist comparison. Weight repacking alone accounted for about 2.5\,GiB of the 3\,B model's CPU memory, which matters on an 8\,GB board. The second is device state: in the pilot, a battery run drifted upward in latency by 15--23\% over about 30 minutes, and idle sleep stalled one request for 17 minutes; during the main grid, 41 minutes on battery slowed the affected configurations until they were re-measured. Neither problem is peculiar to our setup, and both are easy to overlook without per-request logs and a power log. Answers, by contrast, were fully reproducible: re-running \RRetimeConfigs{} configurations reproduced \RRetimeSame{} of \RRetimeTotal{} decisions.

\subsection{Limitations}
\label{sec:limitations}
\textit{Domain gap.} Both datasets are indoor, staged recordings in homes, offices and a laboratory, with clean floors, adequate lighting and few occluders. Industrial floors add machinery, pallets, high-visibility clothing, hard hats and partial occlusion. We did not test on industrial footage, so our figures are best read as an upper bound on factory performance until confirmed there.

\textit{Two datasets, few actors, staged falls.} MCF has a single actor and UR Fall a few; falls are performed, and in MCF mostly onto a mattress, which may itself serve as a cue. Real collapses look different. UR Fall's overhead camera recorded only falls, so overhead false-alarm rates cannot be measured. Some hard-negative types are small (16 kneeling and 22 crouching test samples), as the confidence intervals reveal.

\textit{One device.} All efficiency figures come from one laptop-class system-on-chip, used as a CPU-only and an integrated-GPU operating point. Embedded boards will be slower and have less memory; relative effects of size, precision, resolution and frame count, and the device-independent costs we report, should transfer better than absolute latencies. Power is the package estimate from \texttt{powermetrics}, which excludes memory and other components, so wall-power figures would be higher.

\textit{Single decisions.} We evaluate single frames or short clips. A deployed system would aggregate decisions over time, for example by requiring several consecutive positives, which curbs false alarms at the cost of delay. We do not model this.

\textit{Zero-shot only.} We fine-tune nothing. Fine-tuning would probably raise accuracy, especially for the smallest models, but demands target-domain training data.

\textit{Labels and thresholds.} Labels were reviewed by an AI assistant, then checked by one human, blind on a sample and with labels visible on all frames; \HumanCheckText{} Independent human labelling of every frame would strengthen the benchmark, and subtype labels (crouching versus kneeling, say) are the most uncertain. The calibrated operating point uses development data from the same datasets as the test set, so it is optimistic about how a threshold transfers to a new site.

\textit{Model snapshot and licences.} Small VLMs evolve quickly; our conclusions apply to the versions listed, and the benchmark and scripts are released so newer models can be added. UR Fall restricts use to non-commercial research and MCF states no licence, constraining what we can redistribute and possibly the commercial reuse of derived results.

\subsection{Ethical considerations}
Camera-based monitoring of workers raises privacy concerns. A \pd{} detector is meant to prompt a check by a person, not to monitor productivity or discipline workers, and local processing keeps video on site. We used only public datasets collected by others and recorded no one. In figures we blur faces and show only frames from a dataset whose licence permits it. Any deployment should comply with local law on workplace monitoring and involve workers or their representatives.

\section{Conclusion}
\label{sec:conclusion}
We studied how far vision--language models can be shrunk, quantized and fed with smaller images before they stop detecting a person who is down, and what each configuration costs on edge hardware. \pending{Two or three sentences with the main findings: smallest acceptable configuration per device; quantization and resolution effects; temporal context; hard negatives that remain difficult; how the pose baseline compares.} Two lessons from the pilot are likely to hold beyond this task: short prompts suit small models better than detailed definitions, and the model's probability of answering ``yes'' supports better decisions than its answer. Future work should evaluate on industrial footage, add temporal aggregation of decisions and test whether light fine-tuning closes the gap between the smallest and largest models.

% Back matter covering both IEEE Access and MDPI Sensors requirements.
% IEEE Access: AI-generated text must be disclosed in the Acknowledgment; author
% biographies go after the references in the IEEE Access template.

\section*{Data and Code Availability}
Label files, the benchmark build scripts, inference wrappers, evaluation code and per-run result files (one row per sample and configuration, with the hardware and software environment of each run) are available at \url{https://github.com/mokshx/fusion_vlm}. Frames are not redistributed; the scripts download them from the official sources of UR Fall \cite{kwolek2014urfall} and MCF \cite{auvinet2010mcfd} under their respective terms, and use the OmniFall annotations \cite{schneider2025omnifall}.

\section*{Conflicts of Interest}
The author declares no conflicts of interest.

\section*{Funding}
This research received no external funding.

\section*{Acknowledgment}
The author thanks the creators of the UR Fall and Multiple Cameras Fall datasets and of the OmniFall annotations for making their data available. Claude (Anthropic), a large language model with vision input, was used to write analysis code, to draft portions of the manuscript text and to carry out the visual review of benchmark labels described in Section~\ref{sec:benchmark}. The author directed this work, reviewed and edited all content, checked a blind sample of the labels and went through all labelled frames, and takes full responsibility for the paper.


\bibliographystyle{IEEEtran}
\bibliography{references}

\appendices
\section{Prompts}
\label{app:prompts}
The prompts are fixed; any change gets a new identifier so that results stay comparable.

\noindent\textbf{Prompt A.}
\begin{quote}\small
Is any person fallen or lying on the floor? Answer yes or no.
\end{quote}

\noindent\textbf{Prompt B.}
\begin{quote}\small
You are a workplace safety monitor looking at a camera image. A `person down' is a person whose body is lying on the floor or ground, with the torso roughly horizontal and touching or very close to the floor, for example after a fall or collapse. The following are NOT a person down: standing, walking, bending over, kneeling, crouching or squatting, sitting on the floor or on a chair, lying on a bed, sofa or other furniture, and lying under equipment while visibly doing maintenance work. Is there a person down in this image? Answer with one word: yes or no.
\end{quote}

\noindent\textbf{Prompt B2.}
\begin{quote}\small
A person is `down' only if their torso is lying on the floor or ground. Is any person down? Answer yes or no.
\end{quote}

\noindent\textbf{Multi-frame prefix} (prepended for $n>1$ frames).
\begin{quote}\small
The following $n$ images are consecutive frames from one camera, in time order. Answer about the last frame, using earlier frames as context.
\end{quote}

\section{Inference settings}
\label{app:settings}
llama.cpp server, one instance per configuration:
\begin{quote}\small\raggedright
\texttt{llama-server -m <model.gguf>}\\
\texttt{~~--mmproj <mmproj-F16.gguf>}\\
\texttt{~~-ngl <0|99> -c 4096 -np 1 -cram 0}\\
\texttt{~~--no-webui [-t <threads>] [-dev none]}
\end{quote}
\texttt{-ngl 0 -dev none} keeps all computation on the CPU, including the vision projector; \texttt{-np 1} uses a single slot; \texttt{-cram 0} disables the host-side prompt cache. Requests use the chat-completions endpoint with images as base64 PNG, \texttt{temperature=0}, \texttt{max\_tokens=16}, \texttt{cache\_prompt=false}, \texttt{logprobs=true}, \texttt{top\_logprobs=10} and streaming. $p_\text{yes}$ is computed from the top-10 log-probabilities of the first generated token, summing tokens that equal ``yes'' or ``no'' after lower-casing and stripping whitespace, and renormalising over the two.

{\sloppy Transformers reference: \texttt{AutoModelFor\allowbreak ImageTextToText} in bf16 (fp16 on GPUs without bf16), the model's chat template, greedy decoding, \texttt{max\_new\_tokens=16}; $p_\text{yes}$ from the softmax over the first-step logits.\par}

\section{Label file format}
\label{app:labels}
One CSV row per sample: \texttt{sample\_id}, \texttt{image\_path} (or several, separated by semicolons, for clips), \texttt{label}, \texttt{subtype}, \texttt{viewpoint}, \texttt{dataset}, \texttt{sequence}, \texttt{camera}, \texttt{frame}, \texttt{source\_label} (the source dataset's own annotation), and free-text \texttt{notes}.

\section{Model revisions}
\label{app:revisions}
Hugging Face repository revisions (first 12 characters of the commit hash) of the weights used. GGUF files are from the repositories listed; Transformers runs use the original repositories.

\begin{center}\footnotesize
\begin{tabular}{ll}
\toprule
Repository & Revision \\
\midrule
LiquidAI/LFM2.5-VL-450M-GGUF & \texttt{1abed04b6fe7} \\
LiquidAI/LFM2.5-VL-1.6B-GGUF & \texttt{36fc16bc9513} \\
LiquidAI/LFM2.5-VL-3B-GGUF & \texttt{6f730e9a2c45} \\
unsloth/Qwen3.5-2B-GGUF & \texttt{f6d5376be1ed} \\
unsloth/Qwen3.5-4B-GGUF & \texttt{e87f176479d0} \\
LiquidAI/LFM2.5-VL-450M & \texttt{fc6221ca597f} \\
LiquidAI/LFM2.5-VL-3B & \texttt{35a118d938ce} \\
Qwen/Qwen3.5-2B & \texttt{15852e8c1636} \\
Qwen/Qwen3.5-4B & \texttt{851bf6e806ef} \\
\bottomrule
\end{tabular}
\end{center}


\end{document}
